%-------------------------------------------------------------------------------
\chapter{Introduction}
\label{sec:intro}
%-------------------------------------------------------------------------------

At lowest order in the Standard Model~(SM) electroweak theory~\cite{GLASHOW1961579,SALAM1964168,PhysRevLett.19.1264} the $W$-boson mass, \mW, can be expressed solely as a function of the $Z$-boson mass, $m_Z$, the fine-structure constant, $\alpha$, and the Fermi constant, $G_\text{F}$. Higher-order corrections introduce an additional dependence of the $W$-boson mass on the gauge couplings and the masses of the heavy particles of the SM, such as the top-quark mass, $m_t$, and the Higgs boson mass, $m_H$~\cite{Sirlin:1980nh,Awramik:2003rn}. In extended theories, the loop corrections receive contributions from additional particles and interactions. The consistency of the SM and potential effects of new physics can therefore be probed by comparing the measured values of \mW with the results of global fits to the relevant physical parameters~\cite{deBlas:2021wap, Haller:2018nnx, Erler:1999ug}. The SM fit yields $\mW^\text{SM} = 80355 \pm 6$~\MeV~\cite{deBlas:2021wap, Haller:2018nnx}.
The present experimental situation is characterised by a significant tension between the precise measurement from the CDF Collaboration, $\mW = 80433.5 \pm 9.4$~\MeV~\cite{CDF:2022hxs}, and the average of the LEP~\cite{ALEPH:2013dgf}, D0~\cite{Abazov:2012bv}, ATLAS~\cite{STDM-2014-18} and LHCb~\cite{LHCb:2021bjt} measurements, $\mW = 80369.2 \pm 13.3$~\MeV~\cite{Amoroso:2023pey}. \todo{Need to add ATLAS and CMS here}.


This analysis aims to determine the \mW using dedicated low-mu data-samples. \todo{add more here}. A dedicated calibration of this data-sample is needed and described in Ref.~\cite{lowMuMw:INT1}.


The present analysis uses a profile likelihood~(PLH)~\cite{Rolke:2004mj} fit and performs a simultaneous determination of \mW together with a set of nuisance parameters describing the experimental and modelling uncertainties. \todo{a few more details on the fit method once it is finalized}. The theory used for the fit is \todo{Add some words here also about the PDF used}


This note is organised as follows: \todo{add this plus introduction}

The note is organised as follows: sections~\ref{sec:samples} and~\ref{sec:objects} describe the data and simulated samples, and the object definitions.
Section~\ref{sec:modelingCorr} describes physics corrections used for the theory modelling. The background estimates are described in Section~\ref{sec:bkgs}.
Section~\ref{sec:eventSel} describes the event selection in the \Wboson\ analysis and presents the corresponding event yields and control plots.
The fit methodology is detailed in Section~\ref{sec:fit}.
Section~\ref{sec:sys} discusses the uncertainties. All results are given in section~\ref{sec:results}. A short conclusion is given in section~\ref{sec:conclusion}.